% !PRINTED 622
\weberpage{622}
\begin{center}
\S.\,171.

\gesperrt{Die Primideale im K\"orper $\Omega_m$.}
\end{center}

% [sic: The printed reference 164 should be 169, theorem I and (7), where the prime-power factorisation is established.]
% [Source shorthand: q is associated with sigma^mu, not generally equal to it as an element; the corresponding principal ideals are equal.]
Wir haben im \S.\,164 gesehen, dass die nat\"urliche Primzahl $q$
die $\mu^{\text{te}}$ Potenz einer Primzahl $\sigma$ im K\"orper $\Omega_m$ ist.

Um alle Primideale, die im K\"orper $\Omega_m$ existiren, zu ermitteln, sind also noch s\"ammtliche von $q$ verschiedene nat\"urliche
Primzahlen $p$ in Primfactoren zu zerlegen, und es sind darunter
die nicht associirten auszusuchen.

Die Grundlage f\"ur diese Untersuchung bilden die S\"atze des
\S.\,162, wenn der dort mit $R$ bezeichnete K\"orper durch den
K\"orper der rationalen Zahlen ersetzt wird.

Jede Zahl $\omega$ des K\"orpers $\Omega_m$ geht durch eine der $\mu$ Substitutionen
\[
s_n = (r, r^n)
\]
in eine bestimmte andere Zahl $\omega_n$ \"uber, die gleichfalls in $\Omega_m$
enthalten ist. Ist
\weq{1}{\omega = \omega_1 = a_0 + a_1 r + a_2 r^2 + \cdots + a_{\mu-1} r^{\mu-1},}
so ist
\weq{2}{\omega_n = a_0 + a_1 r^n + a_2 r^{2n} + \cdots + a_{\mu-1} r^{(\mu-1)n}.}
Wenn eine dieser Zahlen $\omega_n$ eine ganze Zahl ist, wie wir
jetzt annehmen wollen, so sind es auch alle anderen, und dies
tritt immer dann und nur dann ein, wenn die rationalen Zahlen
$a_0, a_1, \ldots, a_{\mu-1}$ ganz sind.

Sind $h, k$ irgend zwei Exponenten, so ist
\[
r^h - r^k = r^k(r^{h-k} - 1),
\]
und dies ist nach \S.\,169, (9) mit einer Potenz von $\sigma$ associirt,
also, wenn $\mathfrak{p}$ ein in $p$ aufgehendes Primideal bedeutet, durch $\mathfrak{p}$
nicht theilbar, ausser wenn $h \equiv k \pmod{m}$ ist. Daraus ergiebt
% [sic: The following assertion needs congruence for every integral omega, not for one fixed omega (omega=1 is a counterexample). Choosing omega=r proves the required assertion about the inertia group X, so the unramifiedness conclusion remains valid.]
sich, dass zwei Zahlen $\omega_h, \omega_k$, nur dann nach dem Modul $\mathfrak{p}$ congruent sein k\"onnen, wenn $h \equiv k \pmod{m}$ ist, und dass also die
Gruppe $X$ des \S.\,162, die der Bedingung
\[
\omega\,|\,\chi \equiv \omega \pmod{\mathfrak{p}}
\]
gen\"ugt, nur die \gesperrt{identische Substitution} enth\"alt. Daraus folgt
aber, dass $g = 1$, d.\,h. dass $p$ \gesperrt{nicht durch das Quadrat
eines Primideales theilbar ist.}
